\begin{abstract}
\noindent
%% Abstract
%% Text of abstract
We use a randomized controlled trial (RCT) to examine the causal effect of information provision on household electricity consumption. Households were randomly assigned to one of two conditions. The information condition ($N=454$) received feedback on their electricity consumption via a smartphone app, while the control condition ($N=452$) did not receive feedback via the app. Electricity consumption was tracked for up to 70 weeks after introducing the app (during-treatment) and up to 3 weeks before (pre-treatment). Using data from Google Analytics, we can distinguish between the intent-to-treat group (all app recipients) and those who actively engaged with the app (and thus effectively received the information treatment). Our analysis yields three main results. First, we identify a statistically significant intent-to-treat (ITT) effect, with households in the information condition reducing their electricity consumption by an average of 6\%. Second, the ITT effect underestimates the true savings for households that actively engage with the app. We find that households with higher engagement levels achieve substantially greater reductions in electricity consumption than those with lower engagement. Third, we observe long-term reductions in electricity consumption even as app engagement declines over time. Statistically significant savings persist for up to 30 weeks after treatment. Beyond this period, the effect stabilizes at a lower level of consumption, but is no longer statistically significant.

\end{abstract}
\medskip
\noindent
\textbf{JEL Codes:} C93, D12, Q41 \\
\textbf{Keywords:} Energy consumption, Information provision, Energy conservation, Salience of information, Environmental behavior, Randomized controlled trials